\documentclass{ntu-minimal-thesis}

% ============================================================================
% 論文資料：一般使用者只需要修改這一區
% ============================================================================
\ThesisSetup{
  font-main     = {Songti TC},
  font-sans     = {Heiti TC},
  font-kai      = {Kaiti TC},
  college-zh    = {社會科學院},
  college-en    = {College of Social Sciences},
  department-zh = {經濟學系},
  department-en = {Department of Economics},
  title-zh     = {從三千公里到黃泉：大雄與鄭莊公的尋母方法},
  title-en     = {From Three Thousand Kilometers to the Yellow Springs: How Nobita and Duke Zhuang Found Their Mothers},
  author-zh    = {範例學生},
  author-en    = {Example Student},
  student-id   = {R00000000},
  advisor-zh   = {派大星教授加},
  advisor-en   = {Professor Patrick Star},
  thesis-date  = {2026-06},
  defense-date = {2026-05-15},
  doi          = {10.1234/ntu.thesis},
}

\begin{document}

\makecover
\makeapprovalpage

\begin{acknowledgements}
感謝指導教授派大星教授加博士的耐心指導與鼓勵，
容許我把《左傳》、任意門和一整瓶藥丸放在同一篇論文裡。
感謝潁考叔與哆啦 A 夢；
他們以不同時代的技術證明，
面對一句說得太重的氣話，最好的顧問未必勸人收回，
有時只是替那句話重新安排地理。

最後，感謝所有曾經說過「再也不要見到你」、
五分鐘後卻開始想念對方的人。
願各位都能找到自己的地道、任意門，
或至少一個願意幫忙想辦法的朋友。
\end{acknowledgements}

\begin{abstracten}
This essay compares two stories of a son who angrily renounces his mother
and almost immediately regrets it:
``Mama o Tazunete Sanzen Kiro Jo'' from \emph{Doraemon}
and the account of Duke Zhuang of Zheng in the \emph{Zuo Tradition}.
In the first, a gadget converts Nobita's anger into a journey of three thousand kilometers;
in the second, an oath postpones reunion until the Yellow Springs.
Doraemon and Ying Kaoshu solve these crises in parallel ways.
Neither asks the son to deny his earlier words.
Instead, each turns a figurative impossibility into a literal route:
one with the Anywhere Door and the other with a tunnel reaching groundwater.
The comparison shows that both stories treat space as a mechanism of remorse.
Their endings are equally ironic: reunion restores the mother--son relationship,
but does not necessarily make that relationship any easier.

\keywords{Doraemon,
  Zuo Tradition,
  filial relations,
  oaths,
  space and narrative.}
\end{abstracten}

\begin{abstractzh}
本文比較兩個兒子一氣之下決定不再見母親，
隨即又感到後悔的故事：
《哆啦 A 夢》〈尋母三千里藥丸〉與《左傳》〈鄭伯克段於鄢〉。
前者以道具把大雄的氣話換算成三千公里的路程，
後者則以誓言把鄭莊公與武姜的重逢推遲至「黃泉」。
哆啦 A 夢與潁考叔採取了結構相似的解法：
他們不要求當事人收回前言，
而是把看似不可能的條件改造成可以完成的路線；
一方使用任意門，另一方則挖掘抵達地下泉水的隧道。
兩篇故事因而都把空間變成呈現後悔的機制。
它們的結局也同樣帶有反諷：
母子雖然重新相見，原有的親子關係卻未必因此變得比較容易。

\keywords{哆啦 A 夢、
  左傳、
  母子關係、
  誓言、
  空間敘事。}
\end{abstractzh}

\makethesistableofcontents
\thesismainmatter

\section{兩句說出口就後悔的話}

《哆啦 A 夢》第 21 卷第 13 話〈尋母三千里藥丸〉中，
大雄被媽媽責罵後，賭氣表示討厭媽媽，希望不要再見到她。
小學館的官方書目將日文原題
〈ママをたずねて三千キロじょう〉列為該卷第 13 話，
台灣播映時則使用〈尋母三千里藥丸〉一名
\citep{fujiko1980,litv2019}。

《左傳》〈鄭伯克段於鄢〉也從一句過重的話開始。
鄭莊公平定其弟共叔段的叛亂後，
將參與其事的母親武姜安置於城潁，並發誓：
「不及黃泉，無相見也。」
一個小學生和一位諸侯的處境當然不同，
但兩人都把當下的憤怒說成了永久的母子分離，
也都很快發現，話一旦說出口，後悔並不會自動使它失效。

\section{三千公里：把氣話換算成里程}

哆啦 A 夢拿出的藥丸有一套精確得近乎殘酷的規則：
吃一顆後，必須移動 300 公尺才能再見到媽媽。
大雄第一次試吃，走完 300 公尺找到媽媽，卻又挨了一頓罵；
他一氣之下吞掉整瓶 10,000 顆，
於是必須累積 3,000 公里的距離才能再次相見
\citep{shogakukan2020,litv2019}。

真正的懲罰並不是走路本身，
而是憤怒消退的速度遠比里程累積得快。
大雄沒過多久便開始後悔、想念媽媽，
但藥丸不接受心情變化作為解約條件。
最後，哆啦 A 夢利用任意門，
先把大雄送到 1,500 公里外，再把他送回來；
兩段位移湊足 3,000 公里，
讓大雄不必真的徒步走完三千公里，也沒有違反道具的規則
\citep{fujiko1980}。

\section{黃泉：把誓言改造成地點}

鄭莊公的難題沒有數字，卻有更難取消的政治分量。
他既然公開發誓不到黃泉不與母親相見，
便不能只因後悔而若無其事地接回武姜。
潁考叔得知莊公的悔意後，提出一個近乎法律文字遊戲的方案：
「若闕地及泉，隧而相見，其誰曰不然？」
只要向下挖掘直到看見黃色的地下泉水，
再於隧道中相見，便能說莊公確實到了「黃泉」，
並未背棄自己的誓言。
莊公依計而行，最後與武姜「遂為母子如初」
\citep{zuozhuan}。

\section{任意門與地道}

哆啦 A 夢與潁考叔扮演的是同一類顧問。
他們都沒有直接否定當事人的氣話，
而是接受話語表面的規則，再改造規則所依附的空間。
任意門把三千公里折疊為兩次開門；
地道則把原本指涉死亡世界的「黃泉」，
重新解釋為工程技術可以抵達的地下水。

兩種方法的差別也很重要。
哆啦 A 夢使用未來科技縮短旅程，
潁考叔則透過語義與儀式維持君主誓言的表面效力。
然而兩人都明白：
要修補一句無法直接收回的話，
有時不必改變那句話，而可以改變說話者所在的位置。

\section{「遂為母子如初」之後}

《左傳》以「遂為母子如初」結束這段插曲，
讀來像是一個圓滿結局；但篇章開頭早已說武姜因莊公寤生而
「遂惡之」。
所謂「如初」，恢復的首先是母子的名分與往來，
未必是一段從此毫無嫌隙的親情。

大山版動畫的收尾把同一層反諷變成笑點：
大雄終於抱住想念已久的媽媽，
媽媽卻仍自顧自地說要叫爸爸教訓他 \citep{oyama1981}。
三千公里的位移讓母子重新相見，
卻沒有讓大雄免於原來的家庭秩序。
兩篇故事因此都沒有承諾「重逢之後一切變好」；
它們給出的只是更樸素的結果：
人繞了遠路、鑽過地道，終於回到那段仍然令人又愛又氣的關係。

\section{結語}

〈尋母三千里藥丸〉與〈鄭伯克段於鄢〉相隔兩千多年，
卻共享一種精巧的敘事結構：
憤怒產生禁令，禁令化為空間，後悔召來顧問，
顧問再用一條字面上合規的路徑完成重逢。
一條路經過任意門，另一條路深入黃泉；
它們最終抵達的，都是「如初」。
而「如初」之所以有趣，正因為它既表示團聚，
也表示媽媽很可能還是會繼續罵人。

\printthesisbibliography

\end{document}
